% Options for packages loaded elsewhere
\PassOptionsToPackage{unicode}{hyperref}
\PassOptionsToPackage{hyphens}{url}
\documentclass[
]{article}
\usepackage{xcolor}
\usepackage[margin=1in]{geometry}
\usepackage{amsmath,amssymb}
\setcounter{secnumdepth}{-\maxdimen} % remove section numbering
\usepackage{iftex}
\ifPDFTeX
  \usepackage[T1]{fontenc}
  \usepackage[utf8]{inputenc}
  \usepackage{textcomp} % provide euro and other symbols
\else % if luatex or xetex
  \usepackage{unicode-math} % this also loads fontspec
  \defaultfontfeatures{Scale=MatchLowercase}
  \defaultfontfeatures[\rmfamily]{Ligatures=TeX,Scale=1}
\fi
\usepackage{lmodern}
\ifPDFTeX\else
  % xetex/luatex font selection
\fi
% Use upquote if available, for straight quotes in verbatim environments
\IfFileExists{upquote.sty}{\usepackage{upquote}}{}
\IfFileExists{microtype.sty}{% use microtype if available
  \usepackage[]{microtype}
  \UseMicrotypeSet[protrusion]{basicmath} % disable protrusion for tt fonts
}{}
\makeatletter
\@ifundefined{KOMAClassName}{% if non-KOMA class
  \IfFileExists{parskip.sty}{%
    \usepackage{parskip}
  }{% else
    \setlength{\parindent}{0pt}
    \setlength{\parskip}{6pt plus 2pt minus 1pt}}
}{% if KOMA class
  \KOMAoptions{parskip=half}}
\makeatother
\usepackage{color}
\usepackage{fancyvrb}
\newcommand{\VerbBar}{|}
\newcommand{\VERB}{\Verb[commandchars=\\\{\}]}
\DefineVerbatimEnvironment{Highlighting}{Verbatim}{commandchars=\\\{\}}
% Add ',fontsize=\small' for more characters per line
\usepackage{framed}
\definecolor{shadecolor}{RGB}{248,248,248}
\newenvironment{Shaded}{\begin{snugshade}}{\end{snugshade}}
\newcommand{\AlertTok}[1]{\textcolor[rgb]{0.94,0.16,0.16}{#1}}
\newcommand{\AnnotationTok}[1]{\textcolor[rgb]{0.56,0.35,0.01}{\textbf{\textit{#1}}}}
\newcommand{\AttributeTok}[1]{\textcolor[rgb]{0.13,0.29,0.53}{#1}}
\newcommand{\BaseNTok}[1]{\textcolor[rgb]{0.00,0.00,0.81}{#1}}
\newcommand{\BuiltInTok}[1]{#1}
\newcommand{\CharTok}[1]{\textcolor[rgb]{0.31,0.60,0.02}{#1}}
\newcommand{\CommentTok}[1]{\textcolor[rgb]{0.56,0.35,0.01}{\textit{#1}}}
\newcommand{\CommentVarTok}[1]{\textcolor[rgb]{0.56,0.35,0.01}{\textbf{\textit{#1}}}}
\newcommand{\ConstantTok}[1]{\textcolor[rgb]{0.56,0.35,0.01}{#1}}
\newcommand{\ControlFlowTok}[1]{\textcolor[rgb]{0.13,0.29,0.53}{\textbf{#1}}}
\newcommand{\DataTypeTok}[1]{\textcolor[rgb]{0.13,0.29,0.53}{#1}}
\newcommand{\DecValTok}[1]{\textcolor[rgb]{0.00,0.00,0.81}{#1}}
\newcommand{\DocumentationTok}[1]{\textcolor[rgb]{0.56,0.35,0.01}{\textbf{\textit{#1}}}}
\newcommand{\ErrorTok}[1]{\textcolor[rgb]{0.64,0.00,0.00}{\textbf{#1}}}
\newcommand{\ExtensionTok}[1]{#1}
\newcommand{\FloatTok}[1]{\textcolor[rgb]{0.00,0.00,0.81}{#1}}
\newcommand{\FunctionTok}[1]{\textcolor[rgb]{0.13,0.29,0.53}{\textbf{#1}}}
\newcommand{\ImportTok}[1]{#1}
\newcommand{\InformationTok}[1]{\textcolor[rgb]{0.56,0.35,0.01}{\textbf{\textit{#1}}}}
\newcommand{\KeywordTok}[1]{\textcolor[rgb]{0.13,0.29,0.53}{\textbf{#1}}}
\newcommand{\NormalTok}[1]{#1}
\newcommand{\OperatorTok}[1]{\textcolor[rgb]{0.81,0.36,0.00}{\textbf{#1}}}
\newcommand{\OtherTok}[1]{\textcolor[rgb]{0.56,0.35,0.01}{#1}}
\newcommand{\PreprocessorTok}[1]{\textcolor[rgb]{0.56,0.35,0.01}{\textit{#1}}}
\newcommand{\RegionMarkerTok}[1]{#1}
\newcommand{\SpecialCharTok}[1]{\textcolor[rgb]{0.81,0.36,0.00}{\textbf{#1}}}
\newcommand{\SpecialStringTok}[1]{\textcolor[rgb]{0.31,0.60,0.02}{#1}}
\newcommand{\StringTok}[1]{\textcolor[rgb]{0.31,0.60,0.02}{#1}}
\newcommand{\VariableTok}[1]{\textcolor[rgb]{0.00,0.00,0.00}{#1}}
\newcommand{\VerbatimStringTok}[1]{\textcolor[rgb]{0.31,0.60,0.02}{#1}}
\newcommand{\WarningTok}[1]{\textcolor[rgb]{0.56,0.35,0.01}{\textbf{\textit{#1}}}}
\usepackage{graphicx}
\makeatletter
\newsavebox\pandoc@box
\newcommand*\pandocbounded[1]{% scales image to fit in text height/width
  \sbox\pandoc@box{#1}%
  \Gscale@div\@tempa{\textheight}{\dimexpr\ht\pandoc@box+\dp\pandoc@box\relax}%
  \Gscale@div\@tempb{\linewidth}{\wd\pandoc@box}%
  \ifdim\@tempb\p@<\@tempa\p@\let\@tempa\@tempb\fi% select the smaller of both
  \ifdim\@tempa\p@<\p@\scalebox{\@tempa}{\usebox\pandoc@box}%
  \else\usebox{\pandoc@box}%
  \fi%
}
% Set default figure placement to htbp
\def\fps@figure{htbp}
\makeatother
\setlength{\emergencystretch}{3em} % prevent overfull lines
\providecommand{\tightlist}{%
  \setlength{\itemsep}{0pt}\setlength{\parskip}{0pt}}
\usepackage{bookmark}
\IfFileExists{xurl.sty}{\usepackage{xurl}}{} % add URL line breaks if available
\urlstyle{same}
\hypersetup{
  pdftitle={Bayesian thinking; α control vs decision error},
  hidelinks,
  pdfcreator={LaTeX via pandoc}}

\title{Bayesian thinking; α control vs decision error}
\author{}
\date{\vspace{-2.5em}}

\begin{document}
\maketitle

\textbf{Testing labs} use the main template: Hypothesis → Visualise →
Assumptions → Conduct → Conclude.

\subsection{Learning objectives}\label{learning-objectives}

\begin{itemize}
\tightlist
\item
  Compute a beta-binomial posterior and interpret its credible intervals
  in terms of a decision.
\item
  Contrast frequentist α-control (type I error rate across repeated
  experiments) with Bayesian decision error for a specific experiment.
\item
  Show how prior sensitivity translates into decision sensitivity.
\end{itemize}

\subsection{Prerequisites}\label{prerequisites}

Probability basics and conjugacy.

\subsection{Background}\label{background}

Frequentist testing controls the long-run rate of falsely rejecting a
true null across repeated experiments; the α = 0.05 convention is a
statement about procedure, not about the current dataset. Bayesian
inference instead returns a posterior distribution of the unknown and
turns that into a decision by weighing costs. For a binomial proportion
with a beta prior, the posterior is beta again, so the arithmetic is
clean enough to follow carefully.

The two frameworks answer different questions. ``Given these data and
this model, what is the probability that the treatment is beneficial?''
is a Bayesian question. ``Across many repetitions of this design, how
often would I reject a true null?'' is a frequentist question. In
applied biomedicine, both questions are relevant; knowing which one you
are answering prevents confused write-ups.

Prior sensitivity is a useful discipline regardless of philosophy. A
posterior that shifts materially when the prior is varied across
reasonable alternatives is less trustworthy than one that does not, and
sharing that sensitivity is the Bayesian analogue of transparent
model-building.

\subsection{Setup}\label{setup}

\begin{Shaded}
\begin{Highlighting}[]
\FunctionTok{library}\NormalTok{(tidyverse)}
\FunctionTok{set.seed}\NormalTok{(}\DecValTok{42}\NormalTok{)}
\FunctionTok{theme\_set}\NormalTok{(}\FunctionTok{theme\_minimal}\NormalTok{(}\AttributeTok{base\_size =} \DecValTok{12}\NormalTok{))}
\end{Highlighting}
\end{Shaded}

\subsection{1. Hypothesis}\label{hypothesis}

A small trial observes \emph{k} = 7 successes in \emph{n} = 20 patients
of a new response-rate biomarker. Is the underlying response rate
meaningfully above 0.2?

\subsection{2. Visualise}\label{visualise}

Prior vs posterior for three priors.

\begin{Shaded}
\begin{Highlighting}[]
\NormalTok{theta }\OtherTok{\textless{}{-}} \FunctionTok{seq}\NormalTok{(}\DecValTok{0}\NormalTok{, }\DecValTok{1}\NormalTok{, }\AttributeTok{length.out =} \DecValTok{400}\NormalTok{)}
\NormalTok{priors }\OtherTok{\textless{}{-}} \FunctionTok{tibble}\NormalTok{(}
  \AttributeTok{name =} \FunctionTok{c}\NormalTok{(}\StringTok{"uniform (Beta(1,1))"}\NormalTok{, }\StringTok{"sceptical (Beta(2,8))"}\NormalTok{, }\StringTok{"flat informative (Beta(4,4))"}\NormalTok{),}
  \AttributeTok{a0   =} \FunctionTok{c}\NormalTok{(}\DecValTok{1}\NormalTok{, }\DecValTok{2}\NormalTok{, }\DecValTok{4}\NormalTok{),}
  \AttributeTok{b0   =} \FunctionTok{c}\NormalTok{(}\DecValTok{1}\NormalTok{, }\DecValTok{8}\NormalTok{, }\DecValTok{4}\NormalTok{)}
\NormalTok{)}
\NormalTok{post }\OtherTok{\textless{}{-}}\NormalTok{ priors }\SpecialCharTok{|\textgreater{}}
  \FunctionTok{rowwise}\NormalTok{() }\SpecialCharTok{|\textgreater{}}
  \FunctionTok{mutate}\NormalTok{(}\AttributeTok{data =} \FunctionTok{list}\NormalTok{(}\FunctionTok{tibble}\NormalTok{(}
    \AttributeTok{theta =}\NormalTok{ theta,}
    \AttributeTok{prior =} \FunctionTok{dbeta}\NormalTok{(theta, a0, b0),}
    \AttributeTok{posterior =} \FunctionTok{dbeta}\NormalTok{(theta, a0 }\SpecialCharTok{+}\NormalTok{ k, b0 }\SpecialCharTok{+}\NormalTok{ n }\SpecialCharTok{{-}}\NormalTok{ k)}
\NormalTok{  ))) }\SpecialCharTok{|\textgreater{}}
  \FunctionTok{unnest}\NormalTok{(data) }\SpecialCharTok{|\textgreater{}}
  \FunctionTok{pivot\_longer}\NormalTok{(}\FunctionTok{c}\NormalTok{(prior, posterior), }\AttributeTok{names\_to =} \StringTok{"density"}\NormalTok{)}

\FunctionTok{ggplot}\NormalTok{(post, }\FunctionTok{aes}\NormalTok{(theta, value, }\AttributeTok{colour =}\NormalTok{ density)) }\SpecialCharTok{+}
  \FunctionTok{geom\_line}\NormalTok{() }\SpecialCharTok{+} \FunctionTok{facet\_wrap}\NormalTok{(}\SpecialCharTok{\textasciitilde{}}\NormalTok{ name) }\SpecialCharTok{+}
  \FunctionTok{labs}\NormalTok{(}\AttributeTok{x =} \FunctionTok{expression}\NormalTok{(theta), }\AttributeTok{y =} \StringTok{"density"}\NormalTok{)}
\end{Highlighting}
\end{Shaded}

\subsection{3. Assumptions}\label{assumptions}

Bernoulli trials, exchangeable, conjugate beta prior. With no prior data
on the response rate in this population, we carry all three priors
through and compare decisions.

\subsection{4. Conduct}\label{conduct}

Posterior summaries.

\begin{Shaded}
\begin{Highlighting}[]
  \FunctionTok{rowwise}\NormalTok{() }\SpecialCharTok{|\textgreater{}}
  \FunctionTok{mutate}\NormalTok{(}
    \AttributeTok{post\_mean =}\NormalTok{ (a0 }\SpecialCharTok{+}\NormalTok{ k) }\SpecialCharTok{/}\NormalTok{ (a0 }\SpecialCharTok{+}\NormalTok{ b0 }\SpecialCharTok{+}\NormalTok{ n),}
    \AttributeTok{ci\_lo =} \FunctionTok{qbeta}\NormalTok{(}\FloatTok{0.025}\NormalTok{, a0 }\SpecialCharTok{+}\NormalTok{ k, b0 }\SpecialCharTok{+}\NormalTok{ n }\SpecialCharTok{{-}}\NormalTok{ k),}
    \AttributeTok{ci\_hi =} \FunctionTok{qbeta}\NormalTok{(}\FloatTok{0.975}\NormalTok{, a0 }\SpecialCharTok{+}\NormalTok{ k, b0 }\SpecialCharTok{+}\NormalTok{ n }\SpecialCharTok{{-}}\NormalTok{ k),}
    \AttributeTok{p\_gt\_0\_2 =} \DecValTok{1} \SpecialCharTok{{-}} \FunctionTok{pbeta}\NormalTok{(}\FloatTok{0.2}\NormalTok{, a0 }\SpecialCharTok{+}\NormalTok{ k, b0 }\SpecialCharTok{+}\NormalTok{ n }\SpecialCharTok{{-}}\NormalTok{ k)}
\NormalTok{  ) }\SpecialCharTok{|\textgreater{}}
\NormalTok{  dplyr}\SpecialCharTok{::}\FunctionTok{select}\NormalTok{(name, post\_mean, ci\_lo, ci\_hi, p\_gt\_0\_2)}
\NormalTok{summ}
\end{Highlighting}
\end{Shaded}

A frequentist analogue for comparison.

\begin{Shaded}
\begin{Highlighting}[]
\FunctionTok{c}\NormalTok{(}\AttributeTok{estimate =}\NormalTok{ bt}\SpecialCharTok{$}\NormalTok{estimate, }\AttributeTok{p\_value =}\NormalTok{ bt}\SpecialCharTok{$}\NormalTok{p.value,}
  \AttributeTok{ci\_lo =}\NormalTok{ bt}\SpecialCharTok{$}\NormalTok{conf.int[}\DecValTok{1}\NormalTok{], }\AttributeTok{ci\_hi =}\NormalTok{ bt}\SpecialCharTok{$}\NormalTok{conf.int[}\DecValTok{2}\NormalTok{])}
\end{Highlighting}
\end{Shaded}

\subsection{5. Conclude}\label{conclude}

\begin{quote}
With 7 responders out of 20, the posterior probability that the response
rate exceeds 0.2 is \texttt{round(summ\$p\_gt\_0\_2{[}1{]},\ 2)} under a
flat prior and \texttt{round(summ\$p\_gt\_0\_2{[}2{]},\ 2)} under a
sceptical prior. A one-sided exact binomial test against p₀ = 0.2 gives
\emph{p} = \texttt{signif(bt\$p.value,\ 2)}.
\end{quote}

All three priors point the same direction, but the strength of the
conclusion depends on the prior. For a go/no-go decision, reporting the
posterior probability of exceeding a decision threshold --- along with
its prior-sensitivity --- is more decision-relevant than a
\emph{p}-value.

\subsection{Common pitfalls}\label{common-pitfalls}

\begin{itemize}
\tightlist
\item
  Treating a 95\% credible interval as a 95\% confidence interval; they
  are calibrated for different inferences.
\item
  Choosing a prior that pushes the posterior toward the desired decision
  and not reporting the sensitivity.
\item
  Comparing a one-sided Bayesian tail probability to a two-sided
  \emph{p}-value.
\end{itemize}

\subsection{Further reading}\label{further-reading}

\begin{itemize}
\tightlist
\item
  Spiegelhalter DJ, Abrams KR, Myles JP, \emph{Bayesian Approaches to
  Clinical Trials and Health-Care Evaluation}.
\item
  McElreath R, \emph{Statistical Rethinking}, ch.~2--3.
\end{itemize}

\subsection{Session info}\label{session-info}

\begin{Shaded}
\begin{Highlighting}[]

\end{Highlighting}
\end{Shaded}

\subsection{See also --- chapter index}\label{see-also-chapter-index}

\begin{itemize}
\tightlist
\item
  \href{../index.Rmd}{Methods in this chapter}
\item
  \href{../../../ch00-find-your-method.Rmd}{Find your method (Chapter
  0)}
\end{itemize}

\end{document}
